\documentclass[letterpaper,11pt]{moderncv}
\moderncvstyle{banking}
\usepackage{xpatch}
\moderncvcolor{black}
\AtBeginDocument{\definecolor{color2}{rgb}{0,0,0}}
\usepackage[letterpaper, margin=0.6in]{geometry}
\usepackage{xcolor}
\usepackage{graphicx}
\usepackage{fontspec}
\setmainfont{Times New Roman}
\pagenumbering{arabic}
\usepackage[unicode]{hyperref}
\sloppy
\usepackage{enumitem}

\usepackage[english]{babel}
\usepackage[autostyle, english = american]{csquotes}
\MakeOuterQuote{"}

%==================== PERSONAL INFO ====================%
\name{\textcolor{black}{Saeyoung}}{Kim}
\address{Cambridge, MA, USA}
\phone[mobile]{+1 (351) 322 5510}
\email{saeyoung\_kim@uml.edu}
\social[linkedin]{saeyoung-macx-kim}
\social[github]{PrimaryAnomaly}

\begin{document}

\makecvtitle
\vspace{-3em}

%==================== SUMMARY ====================%
\section{Professional Summary}
Mechatronics engineer with 9+ years of experience designing, building, and validating complex electromechanical systems from concept through production deployment. Designed and integrated sensor suites, fluidic actuator systems, embedded control electronics, and custom ground support equipment for satellite and space payload programs. Deep hands-on experience in mechanism design, PCB development, firmware integration, and system-level commissioning. Proven track record of structured root cause analysis, DFM application, and cross-functional coordination in high-reliability environments.

\section{Technical Skills}
\begin{minipage}[t]{0.48\textwidth}
\cvitem{Mechanical}{Mechanism Design, Tolerance Analysis, Pneumatics\slash Fluidics, Thermal Management, GD\&T}
\cvitem{Electrical}{Circuit Design, PCB Layout, Power Distribution, Signal Integrity, Sensor Integration}
\cvitem{Firmware}{C/C++ (MCU), Real-Time Control, Data Acquisition, Autonomous Operation}
\end{minipage}%
\hfill
\begin{minipage}[t]{0.48\textwidth}
\cvitem{Software}{Python, MATLAB, Git, Linux}
\cvitem{CAD/Fabrication}{Fusion360, AutoCAD, KiCad, 3D Printing, CNC, Laser Cutting}
\cvitem{Methods}{DFM\slash DFA, System Validation, Root Cause Analysis, V\&V, SPC}
\end{minipage}

%==================== EXPERIENCE ====================%
\section{Professional Experience}

\cventry{2025.07--Present}{\textbf{Postdoctoral Associate} | Electromechanical System Automation and Process Control}{University of Massachusetts Lowell}{Lowell, MA, USA}{}{
\begin{itemize}[leftmargin=1cm, itemsep=0pt, parsep=0pt]
    \item Designed and deployed automated data acquisition and control systems for multi-day manufacturing process cycles
    \item Developed Python-based automation and diagnostics tooling, replacing commercial platforms with scalable in-house solutions
    \item Performed system validation and structured root cause analysis on process equipment failures
    \item Architected autonomous process optimization platform coupling software agents with physics-based simulation models
    \item Coordinated cross-functional deliverables across engineering and research teams for government-funded collaboration
\end{itemize}}

\cventry{2024.02--2025.05}{\textbf{Systems Integration Engineer} | Electromechanical Integration and Commissioning}{Hanwha Systems}{Seoul, South Korea}{}{
\begin{itemize}[leftmargin=1cm, itemsep=0pt, parsep=0pt]
    \item Designed and commissioned Electrical Ground Support Equipment (EGSE) and automation cell configurations for satellite assembly and integration
    \item Led electromechanical integration of military SAR satellite systems from subsystem build through system-level acceptance
    \item Selected and integrated electrical, mechanical, and sensor components across power, communications, and control subsystems
    \item Performed system validation, troubleshooting, and structured root cause analysis on integration anomalies
    \item Provided DFM and DFA feedback to design teams, influencing subsystem design for ease of assembly and production scaling
    \item Developed integration procedures, acceptance criteria, and commissioning documentation for complex hardware systems
    \item Contributed to factory architecture design for a high-throughput satellite manufacturing facility
\end{itemize}}

%==================== EDUCATION ====================%
\section{Education}
\cventry{2017.03--2024.02}{\textbf{PhD in Electrical Engineering} | Electromechanical Payload Systems}{Korea University}{Seoul, South Korea}{}{
\begin{itemize}[leftmargin=1cm, itemsep=0pt, parsep=0pt]
    \item Designed and built multi-modal sensor suites integrating electrochemical, optical, and thermal sensors with custom actuator mechanisms
    \item Developed embedded firmware (C/C++) for real-time control, data acquisition, and autonomous system operation on MCU platforms
    \item Designed and fabricated custom PCBs and mechanical assemblies, applying DFM and tolerance analysis from prototype through validation
    \item Built and commissioned pressurized fluidic and thermal management subsystems with pneumatic and electromechanical actuation
    \item Performed system-level validation and environmental testing including thermal cycling and mechanical stress characterization
    \item Secured and managed \$220K research grant; trained junior researchers on build procedures and equipment operation
    \item Published 2 first-author and 6 co-authored journal papers on embedded sensor systems and space payloads
\end{itemize}}

\cventry{2013.03--2017.02}{\textbf{BSc in Mechanical Engineering}}{Konkuk University}{Seoul, South Korea}{}{
\begin{itemize}[leftmargin=1cm, itemsep=0pt, parsep=0pt]
    \item Conducted research in sensor fabrication, electromechanical prototyping, and materials characterization
    \item Gained experience in CAD modeling, GD\&T, CNC machining, PCB prototyping, and precision measurement
\end{itemize}}

%==================== PUBLICATIONS ====================%
\section{Selected Publications}
\begin{sloppypar}
\begin{itemize}[leftmargin=1cm, itemsep=0pt, parsep=0pt]
    \item Jinhwee Kil, \textbf{\underline{Saeyoung Kim}}, James Jungho Pak, "Investigation of thermal management strategies for biological CubeSat payloads," Acta Astronautica, 2024, Art. no. 228.
    \item \textbf{\underline{Kim, S.}}; Park, S.; Pak, J.J., "Multi-Modal Multi-Array Electrochemical and Optical Sensor Suite for a Biological CubeSat Payload," MDPI Sensors, 2024, 24, Art. no. 265.
    \item \textbf{\underline{Saeyoung Kim}}, Ji-Hoon Han, Beelee Chua, James Jungho Pak, "A pH Sensing Pipette for Cross-Contamination Prevention in Industrial Fermentation," IEEE Trans. Industrial Electronics, 2022, 69(7), 7461-7469.
\end{itemize}
\end{sloppypar}

\vfill
\centering{References available upon request.}

\end{document}
